%% kapa: exam=1 round=31 year=2020 period=1 form=A subject=economics subject_ko=경제학원론 seq=33 number=73
%% source: Q-Net 감정평가사 기출문제 (cst00309) artlSeq=5207219 file=extracted/1008/5207219/A형/A형/2020년도 제31회 감정평가사 1차 시험_1교시_A형.pdf
%% conversion: kordoc 4.24.51 -> tools/gen1.py (auto); review=proofread-sonnet; flags=none
\begin{question}{73}
자본 이동이 완전한 먼델-플레밍(Mundell-Fleming)모형에서 A국의 정부지출 확대 정책의 효과에 관한 설명으로 옳은 것은? (단, A국은 소규모 개방경제이며, A국 및 해외 물가 수준은 불변, IS 곡선은 우하향, LM 곡선은 우상향)
\begin{choices}
  \item 환율제도와 무관하게 A국의 이자율이 하락한다.
  \item 고정환율제도에서는 A국의 국민소득이 증가한다.
  \item 변동환율제도에서는 A국의 국민소득이 감소한다.
  \item 고정환율제도에서는 A국의 경상수지가 개선된다.
  \item 변동환율제도에서는 A국의 통화가치가 하락한다.
\end{choices}
\end{question}
\answer{2}
